\section{Where People Remain in the Loop}

The useful place for a person cannot be read from a permanent list of uniquely human skills. Machines change. Tasks change. Institutions assign rights and responsibilities. Three questions make the comparison more durable: who performs the whole task efficiently, who improves the feedback structure, and who has standing to decide that a result is good enough?

\subsection{The Whole Energy Bill}

The adult brain operates at a power level on the order of tens of watts, while the whole body has a larger, activity-dependent metabolic demand \cite{smil2015power}. Comparing that brain figure with the peak power of a data center mixes boundaries.

For system $s$, a useful decomposition is
\begin{equation}
 E_s=E_{s,\mathrm{sensing+computation}}
     +E_{s,\mathrm{execution}}
     +E_{s,\mathrm{support}},
 \quad s\in\{\mathrm{human},\mathrm{machine}\}.
\end{equation}
All terms use the same energy unit, task, quality standard, and completion window. Categories must be mutually exclusive. Where biological functions cannot be separated cleanly, a total measured at the chosen boundary is better than a fictional allocation.

A maintenance worker in an unfamiliar building handles stairs, labels, obstructed access, tools, conversation, and an incomplete description of the fault. A machine serving the same role may need remote support, sensing redundancy, and recovery equipment. The person's training, fatigue, injuries, transport, and organizational support also belong in the chosen comparison.

Body mechanics can reduce the need for active correction on both sides. So can environment design. Standardized racks and containers change the task, often enough to make automation attractive. An open environment and a human-shaped environment are not identical categories: many supposedly natural tasks already benefit from years of design around people.

Task-matched measurement determines the energy result. Reliability, cost, dignity, and service quality then determine how much that result should matter.

\subsection{Different Responses Can Stabilize a System}

Uniform automation can create common responses. Many participants read the same signal, use similar models, and act together. A local correction becomes a collective shock.

For linear channels responding at one frequency, the combined amplitude is
\begin{equation}
 G_{\mathrm{sum}}=
 \left|\sum_{i=1}^{N}g_i e^{\mathrm{i}\varphi_i}\right|.
\end{equation}
Here $g_i$ is an amplitude gain and $\varphi_i$ a phase in radians. Aligned phases add coherently; dispersed phases can partly cancel. This is an amplitude relation. Closed-loop stability requires the complete feedback structure.

Randomized backoff, staggered execution, independent review, and differing information sources can reduce harmful synchronization. People sometimes provide such differences. They can also share biases, follow the same incentives, or respond too late.

The goal is useful diversity of response. It can be supplied by people, software, institutions, or a combination. Its value has to appear in the system's behavior.

Human supervision has particular requirements. The operator needs relevant information, time to act, maintained skill, and real authority. Research on the ironies of automation and human use of automation explains why a nominal person in the loop can be a weak safeguard \cite{bainbridge1983ironies,parasuraman1997humans}. A person asked to rescue an opaque process after long inactivity may have less usable control than the diagram suggests.

\subsection{Meaning Enters the Stopping Rule}

A search can stop because it has reached a time budget, an error threshold, or a convergence test. A society also needs to decide when the result may be relied upon.

Let $\mathcal A$ be the set of states meeting a specified quality and risk requirement. A first arrival can be written
\begin{equation}
 t^\ast=\inf\{t\geq0:\mathbf{x}_t\in\mathcal A\},
 \qquad \inf\varnothing=+\infty.
\end{equation}
The expression defines a stopping event. A real program can also exhaust its budget and report failure. The important social content lies in choosing and revising $\mathcal A$.

Which errors are tolerable? Whose risk counts? What evidence permits deployment? Who can contest the decision? These questions bring meaning into the admission conditions for action. They can be informed by models, yet their authority comes from the people and institutions entitled to set them.

Metrics become part of behavior when attached to rewards and penalties. Goodhart's and Campbell's observations describe why a successful measurement can become a target that distorts the activity being measured \cite{goodhart1975problems,campbell1979assessing}. Admission conditions need review and appeal as well as automatic checking.

Expertise and standing are different. A specialist may best estimate a failure probability. The people exposed to that failure still have interests and rights in deciding whether the exposure is acceptable. A workable system connects the two through intelligible evidence, accountable decisions, and remedies.

The human position is therefore specific. Energy claims need fair measurement. Stabilizing roles need a demonstrated feedback benefit. Decision rights need legitimate authorization. As technology changes, these criteria can guide a changing division of work without turning every capability update into a referendum on human worth.
